\section*{第 \ref{ch:b2:phylogenetic-trees} 章——系统发生树与分子钟}

\begin{solution}{exo:b2:phylogenetic-trees:1}
\omterm{def:b2:phylogenetic-trees:tree}{同源性}：由共同祖先遗传而来的性状——
蝙蝠翅膀与鲸鳍中的肱骨、桡骨和尺骨。
\omterm{def:b2:phylogenetic-trees:tree}{同功性}：独立演化出的相似性状——蝙蝠与鸟类作为飞行面的翅膀
（作为前肢，这些骨骼是同源的，翅膀本身则不是），
或章鱼与脊椎动物的照相机式眼睛。
\omterm{def:b2:phylogenetic-trees:tree}{共有衍征}：由某一\omterm{def:b2:phylogenetic-trees:tree}{演化支}共有并借以界定该\omterm{def:b2:phylogenetic-trees:tree}{演化支}的衍生\omterm{def:b2:phylogenetic-trees:tree}{同源性}——
鸟类的羽毛，
羊膜动物的羊膜卵。
\end{solution}

\begin{solution}{exo:b2:phylogenetic-trees:2}
鸟类：单系。不含鸟类的爬行动物：并系。
不含四足动物的鱼类：并系（肺鱼与我们的
亲缘比与鳟鱼更近）。会飞的脊椎动物（蝙蝠、鸟类、翼龙）：多系。
哺乳动物：单系。\omterm{def:b2:unicellular-diversity:prokaryote}{原核生物}：并系（细菌和古菌，
不含真核生物——而古菌与真核生物的亲缘
比与细菌更近）。
\end{solution}

\begin{solution}{exo:b2:phylogenetic-trees:3}
鸟类的姐妹群：鳄类。哺乳动物的姐妹群：
其余所有羊膜动物（龟、蜥蜴、鳄类、鸟类合在一起）。包含蜥蜴与鸟类的最小
\omterm{def:b2:phylogenetic-trees:tree}{演化支}：（蜥蜴，（鳄类，鸟类）），
即本树中不含龟的双孔类。在其节点处交换鳄类与
鸟类，不改变任何结论：节点可以自由旋转。
\end{solution}

\begin{solution}{exo:b2:phylogenetic-trees:4}
无根树：$3$；$3\times 5\times 7 = 105$；$3\times 5\times
7\times 9\times 11\times 13\times 15 = \num{2027025}$。4 个分类单元的有根树
数为 15：3 棵无根树各有 5 条分支，共得 15 种。
\end{solution}

\begin{solution}{exo:b2:phylogenetic-trees:5}
$((A,B),(C,D))$：位点 1、2、3 各 1 步，位点 4 为 2 步
（两侧分别为 $\mathrm{A}$ 与 $\mathrm{G}$），位点 5 为 1 步，共 6 步。
$((A,C),(B,D))$：$2 + 2 + 2 + 1 + 1 = 8$。$((A,D),(B,C))$：$2 + 2 + 2 +
2 + 1 = 9$。第一棵树最简约；在其上，位点 4 是同塑的（$\mathrm{A}\to\mathrm{G}$
的变化，或其逆变化，
发生两次）。
\end{solution}

\begin{solution}{exo:b2:phylogenetic-trees:6}
在高度 3 处将 $A$ 与 $B$ 连接。于是 $d(AB,C) = 14$，$d(AB,D) = 16$，
$d(C,D) = 10$：在高度 5 处将 $C$ 与 $D$ 连接。于是 $d(AB,CD) = (14 + 16
+ 14 + 16)/4 = 15$：根置于高度 $7.5$。树为 $((A,B),(C,D))$。
\end{solution}

\begin{solution}{exo:b2:phylogenetic-trees:7}
$d = -0.75\ln(1 - 4p/3)$：$0.052$、$0.38$、$1.21$。当 $p = 0.6$ 时，
对数的真数为 $0.2$，斜率 $\mathrm{d}d/\mathrm{d}p
= 1/(1 - 4p/3) = 5$：抽样误差 $0.02$（针对 $p$）
变为 $0.1$（在 $d$ 中），而假设每个位点速率相等
（这在任何真实基因中都不成立）引入同样大小的偏倚。该基因
已饱和。
\end{solution}

\begin{solution}{exo:b2:phylogenetic-trees:8}
$t = d/2k = 0.02/(2\times 10^{-9}) = 10$ 百万年。在 500
百万年时，$d = 2\times 10^{-9}\times 5\times 10^{8} = 1.0$，且
$p = 0.75(1 - e^{-4/3}) = 0.55$：距饱和已走过四分之三，
此时校正急剧变陡，年代不可靠，
——对这类\omterm{def:b2:cell-differentiation:differentiation}{分化}，需要使用进化更慢的基因。
\end{solution}

\begin{solution}{exo:b2:phylogenetic-trees:9}
\qty{52}{\%}：该\omterm{def:b2:phylogenetic-trees:tree}{演化支}仅在刚过半数的重抽样中得以恢复；
数据对其几乎没有信息，
各替代方案合起来的可能性与之相当。\qty{99}{\%}：
各位点之间信号一致，该\omterm{def:b2:phylogenetic-trees:tree}{演化支}得到强有力的支持
（以所用模型和比对为前提——\omterm{prop:b2:phylogenetic-trees:pitfalls}{长枝吸引}等系统误差
无法被自展检验探测到）。对于前一种情形，
可增加序列（更多基因）、增加能切分长枝的分类单元，
并改用另一种方法检验。
\end{solution}

\begin{solution}{exo:b2:phylogenetic-trees:10}
两个谱系各自在很大比例的位点上发生了改变，
在两者都改变的任一位点上，落到同一碱基的概率为四分之一
（设四种碱基速率相同）。若二者各自在
\qty{40}{\%} 的位点上发生改变，则二者同时改变的位点为
\qty{16}{\%}，偶然一致的位点占全部位点的 \qty{4}{\%}——
这一比例可能超过携带真实\omterm{def:b2:phylogenetic-trees:tree}{共有衍征}的位点比例，
而这些\omterm{def:b2:phylogenetic-trees:tree}{共有衍征}本可把各自与其真实的、进化缓慢的
近缘类群联系起来。\omterm{meth:b2:phylogenetic-trees:parsimony}{简约法}只计数相同之处，不追究其成因，
便把二者连在一起。在每条长枝上增加分叉出的分类单元，
可把长枝切分为较短的片段，使偶然的相同之处
分散到若干较短的分支上，
每一片段的真实\omterm{def:b2:phylogenetic-trees:tree}{共有衍征}便得以显现；而似然模型
若预期长枝上会有大量偶然相同之处，则
可直接对其加以校正。
\end{solution}

\begin{solution}{exo:b2:phylogenetic-trees:11}
$d = -0.75\ln(1 - 0.12) = 0.096$; $t = d/2k = 0.096/(4\times 10^{-8})
= 2.4$ 百万年。替换数：$0.096\times \num{16000} =
1500$，相对精度 $1/\sqrt{1500} = \qty{2.6}{\%}$：统计误差
为 $\pm 0.06$ 百万年。但该速率是依据其他\omterm{def:b2:cell-differentiation:differentiation}{分化}校准的，对本谱系
可能相差两倍，
原因在于谱系间速率存在差异，以及世代尺度上测得的系谱速率
与数百万年尺度上测得的速率并不相同；
此外，该基因的分歧早于物种的\omterm{def:b2:cell-differentiation:differentiation}{分化}，
因为祖先种群存在多态性，
故真正的\omterm{def:b2:cell-differentiation:differentiation}{分化}时间比基因所示的更近，
其差距取决于祖先种群的大小。
（由大量核基因与化石得出的公认年代为 600 万至 700 万年，
此处的线粒体分子钟
偏快。）
\end{solution}

\begin{solution}{exo:b2:phylogenetic-trees:12}
\omterm{prop:b2:phylogenetic-trees:pitfalls}{不完全谱系分选}：当两次物种形成在
时间上相近时，祖先多态性可能发生分选，使基因树
对物种的分组与物种树不同；水平转移：
从另一谱系获得的基因携带该谱系的
历史；基因复制：比较一个物种中的旁系同源基因与
另一物种中的直系同源基因，得到的是复制的年代，
而非物种形成的年代。物种树的获得方法是：利用
来自全基因组的许多基因，取多数支持的树
（或采用一种预期已知比例不一致基因的模型），
通过谨慎选择直系同源基因，并在
\omterm{def:b2:unicellular-diversity:prokaryote}{原核生物}中优先选用
很少被转移的核糖体基因和信息类基因。
\end{solution}

\begin{solution}{pb:b2:phylogenetic-trees:1}
\textbf{1.} 位点 1、2 和 4（仅由 $W$ 与 $H$ 共有的状态，
相对于骆驼为衍生状态）：$\{W,H\}$ 的\omterm{def:b2:phylogenetic-trees:tree}{共有衍征}。位点 3
和 6：$\{W,H,C\}$ 的\omterm{def:b2:phylogenetic-trees:tree}{共有衍征}。位点 5 是 $W$ 的独有衍征，
即在单一谱系中发生的变化，不构成任何类群；在分子树上没有位点
是同塑的。
\textbf{2.} 每个位点一次变化：共 6 步。
\textbf{3.} 位点 1、2 和 4 各需两次变化，位点 3、5
和 6 各一次：共 9 步。
\textbf{4.} 鲸–河马树领先三步。
在一千个位点中仅差三个，这一微小优势可能被少数
同塑性推翻；要加强该结论，需要更多基因和更多
分类单元（其他反刍动物、美洲野猪、鹿），并进行自展检验。
\textbf{5.} 该\omterm{def:b2:phylogenetic-trees:tree}{演化支}在六次重抽样中有五次重现：支持度中等，
高于机遇水平，但未达到通常所称强支持的
\qty{95}{\%}。
\textbf{6.} \omterm{meth:b2:phylogenetic-trees:parsimony}{外群}必须位于类群之外，但需足够接近，
以使其序列仍可比对且未饱和；骆驼是
位于牛–猪–河马–鲸类群之外的偶蹄类动物。若以
鱼类为\omterm{meth:b2:phylogenetic-trees:parsimony}{外群}，其差异已接近饱和水平，其状态
无法说明哺乳动物内部哪种状态是祖先的，
且其长枝会吸引进化最快的内群
谱系。
\textbf{7.} $d = 0.0625$，$0.107$，$0.131$，$0.155$。
\textbf{8.} 连接 $W$ 与 $H$，高度为 $0.031$。然后 $d(WH,C) =
0.107$，$d(WH,P) = 0.131$，$d(C,P) = 0.131$，而到 $M$ 的距离为
$0.155$：连接 $WH$ 与 $C$，高度为 $0.054$。然后 $d(WHC,P) = 0.131$：
在 $0.065$ 处连接。最后连接 $M$，高度为 $0.078$。树为 $(M,(P,(C,(W,H))))$。
\textbf{9.} 是：树相同，嵌套关系也相同。
\textbf{10.} 若鲸的变化快了一倍，它与所有类群的距离都会更远，
包括河马在内；\omterm{meth:b2:phylogenetic-trees:upgma}{UPGMA} 方法将每个末端置于高度零，
并连接最近的一对，会先把河马与牛连接，
将鲸置于外侧——恰好得到解剖学家的树，
但理由是错误的。
\textbf{11.} $0.0625\times 1000 = 62.5$ 次替换；泊松
标准差 $\sqrt{62.5} = 7.9$。
\textbf{12.} $7.9/62.5 = \qty{13}{\%}$。
\textbf{13.} $k = d/2t = 0.107/(1.2\times 10^{8}) = 8.9\times 10^{-10}$
每位点每年。
\textbf{14.} $t = 0.0625/(2\times 8.9\times 10^{-10}) = 35$ 百万
年。
\textbf{15.} 猪：$0.131/(1.79\times 10^{-9}) = 73$ 百万年；
骆驼：$0.155/(1.79\times 10^{-9}) = 87$ 百万年。
\textbf{16.} $35 \pm 4.5$ 百万年（\qty{13}{\%}）。
\textbf{17.} 速率正比于 $1/t_{\text{cal}}$，年代正比于
$t_{\text{cal}}$：$\pm 5/60 = \qty{8}{\%}$，因此 $\pm 3$ 百万年
——与泊松误差合并后，约为 $\pm 5$ 百万年。
\textbf{18.} $d = -0.75\ln(1 - 0.6) = 0.69$：该基因实际上已在
约三分之二的位点发生变化，校正使
$p$ 放大 $1.5$ 倍，斜率为 $2.5$；该分歧已接近饱和，
其年代参考价值甚微。
\textbf{19.} 某一性状的缺失并非共有的衍生状态：
鲸完全失去了后肢及其踝部，因此无从
推断它们原本具有何种踝部结构。\omterm{def:b2:phylogenetic-trees:tree}{共有衍征}
是\omterm{def:b2:phylogenetic-trees:tree}{演化支}的证据；它在某一谱系中的丢失
并不构成其不属于该\omterm{def:b2:phylogenetic-trees:tree}{演化支}的证据。
\textbf{20.} 化石鲸具有双滑轮踝：鲸
源自具有该结构的祖先，这使它们位于
偶蹄类之中。分子预测得到
岩层记录的证实。
\textbf{21.} 并系：一个祖先连同其全部后代
但不包括鲸谱系。
\textbf{22.} 祖先多态性的\omterm{prop:b2:phylogenetic-trees:pitfalls}{不完全谱系分选}，发生于两次快速
物种形成之间（牛、河马和鲸在数百万年内\omterm{def:b2:cell-differentiation:differentiation}{分化}），或
比较旁系同源基因。判定方法是
对大量基因测序，取多数基因以及串联
分析所支持的树；并检查各\omterm{prop:b2:genome-diversification:duplication}{基因家族}中是否存在重复。
\textbf{23.} $d = -0.75\ln(1 - 0.333) = 0.30$，$k = d/2t =
0.30/(7\times 10^{7}) = 4.3\times 10^{-9}$，每位点每年，
为核基因的五倍。
\textbf{24.} 线粒体 DNA 由校对能力较弱的聚合酶复制，
并暴露于呼吸链产生的氧自由基之中，
且不能由核内修复机制修复；它
也缺乏重组，因此损伤不断累积。
\textbf{25.} 树 $(M,(P,(C,(W,H))))$：鲸是河马的姐妹
群；$k = 8.9\times 10^{-10}$，每位点每年；鲸–河马
\omterm{def:b2:cell-differentiation:differentiation}{分化}时间为 $35 \pm 5$ 百万年。
\end{solution}
